% Options for packages loaded elsewhere
\PassOptionsToPackage{unicode}{hyperref}
\PassOptionsToPackage{hyphens}{url}
\PassOptionsToPackage{dvipsnames,svgnames,x11names}{xcolor}
%
\documentclass[
  11pt,
]{article}
\usepackage{amsmath,amssymb}
\usepackage{iftex}
\ifPDFTeX
  \usepackage[T1]{fontenc}
  \usepackage[utf8]{inputenc}
  \usepackage{textcomp} % provide euro and other symbols
\else % if luatex or xetex
  \usepackage{unicode-math} % this also loads fontspec
  \defaultfontfeatures{Scale=MatchLowercase}
  \defaultfontfeatures[\rmfamily]{Ligatures=TeX,Scale=1}
\fi
\usepackage{lmodern}
\ifPDFTeX\else
  % xetex/luatex font selection
\fi
% Use upquote if available, for straight quotes in verbatim environments
\IfFileExists{upquote.sty}{\usepackage{upquote}}{}
\IfFileExists{microtype.sty}{% use microtype if available
  \usepackage[]{microtype}
  \UseMicrotypeSet[protrusion]{basicmath} % disable protrusion for tt fonts
}{}
\makeatletter
\@ifundefined{KOMAClassName}{% if non-KOMA class
  \IfFileExists{parskip.sty}{%
    \usepackage{parskip}
  }{% else
    \setlength{\parindent}{0pt}
    \setlength{\parskip}{6pt plus 2pt minus 1pt}}
}{% if KOMA class
  \KOMAoptions{parskip=half}}
\makeatother
\usepackage{xcolor}
\usepackage[margin=1in]{geometry}
\usepackage{longtable,booktabs,array}
\usepackage{calc} % for calculating minipage widths
% Correct order of tables after \paragraph or \subparagraph
\usepackage{etoolbox}
\makeatletter
\patchcmd\longtable{\par}{\if@noskipsec\mbox{}\fi\par}{}{}
\makeatother
% Allow footnotes in longtable head/foot
\IfFileExists{footnotehyper.sty}{\usepackage{footnotehyper}}{\usepackage{footnote}}
\makesavenoteenv{longtable}
\setlength{\emergencystretch}{3em} % prevent overfull lines
\providecommand{\tightlist}{%
  \setlength{\itemsep}{0pt}\setlength{\parskip}{0pt}}
\setcounter{secnumdepth}{-\maxdimen} % remove section numbering
\ifLuaTeX
  \usepackage{selnolig}  % disable illegal ligatures
\fi
\IfFileExists{bookmark.sty}{\usepackage{bookmark}}{\usepackage{hyperref}}
\IfFileExists{xurl.sty}{\usepackage{xurl}}{} % add URL line breaks if available
\urlstyle{same}
\hypersetup{
  pdfauthor={Stefano Alimonti --- stefalimonti@gmail.com},
  colorlinks=true,
  linkcolor={blue},
  filecolor={Maroon},
  citecolor={Blue},
  urlcolor={blue},
  pdfcreator={LaTeX via pandoc}}

\author{Stefano Alimonti --- stefalimonti@gmail.com}
\date{March 2026}

\begin{document}

{
\hypersetup{linkcolor=black}
\setcounter{tocdepth}{2}
\tableofcontents
}
\section{Carry Polynomials and the Euler Product: An Approximation
Framework}\label{carry-polynomials-and-the-euler-product-an-approximation-framework}

\textbf{Author:} Stefano Alimonti \textbf{Affiliation:} Independent Researcher
\textbf{Date:} March 2026

\begin{center}\rule{0.5\linewidth}{0.5pt}\end{center}

\newpage

\subsection{Abstract}\label{abstract}

We show that positional integer multiplication in base \(b\) generates companion
matrices \(M_l\) --- one per prime \(l\) --- whose ensemble-averaged spectral
determinant approximates individual Euler factors of the Riemann zeta function:

\[\langle|\det(I - M_l/l^s)|\rangle_N = \frac{R(l,s)}{|1 - l^{-s}|}\]

where \(R(l,s) = 1 + O(l^{-\sigma})\) is a genuine correction factor
(\(R \neq 1\)). Two structural results are proven rigorously: the \textbf{Carry
Representation Theorem} (CRT: the quotient polynomial coefficients equal negated
carries, \(q_i = -c_{i+1}\)) and the \textbf{Unit Leading Carry} theorem (ULC:
\(c_{\text{top}} = 1\) in base 2). The correction factor admits the universal
expansion \(\ln R = c_1/l^s + c_2/l^{2s} + O(l^{-3\sigma})\), where
\(c_1 = \langle c_{\text{top}-1}\rangle - 1\) is the \textbf{trace anomaly}.
High-statistics Monte Carlo (\(\geq 10^8\) samples at \(d = 32\)--\(128\))
establishes \(c_1 = 0.17460 \pm 0.00003\), ruling out \(\ln(2)/4\) at
\(> 30\sigma\). Exact enumeration up to \(K = 21\)
(\(4^{21} \approx 4.4 \times 10^{12}\) configurations) confirms the convergence
of \(c_1(K)\) toward \(\pi/18 = 0.17453\ldots\) (matching to 4.3 significant
digits at \(K = 21\)); the residual is consistent with the polynomial-geometric
convergence model \(c_1(K) = \pi/18 + P(K)(1/2)^K\) {[}E{]}. The appearance of
\(\pi\) is explained by the Angular Uniqueness Theorem: base 2 is the unique
base for which the D-parity boundary is a straight line
(\(\alpha + \beta = \pi/4\)) in angular coordinates, pulling \(\pi\) from the
integration limit {[}E{]}. For \(b > 2\), \(c_1(b)\) does not involve \(\pi\).
The framework provides a concrete arithmetic mechanism linking positional
multiplication to \(\zeta(s)\) as an approximation, not an exact identity.

\textbf{Keywords:} Riemann zeta function, Euler product, carry polynomials,
companion matrices, trace anomaly

\textbf{MSC 2020:} 11M06, 11A63, 15A18, 60J10

\begin{center}\rule{0.5\linewidth}{0.5pt}\end{center}

\newpage

\subsection{1. Introduction}\label{introduction}

The Riemann zeta function \(\zeta(s) = \prod_p(1-p^{-s})^{-1}\) encodes the
distribution of primes through its Euler product. We show that this product
arises, approximately, from carries in positional multiplication.

When two integers \(p, q\) are multiplied in base \(b\), the carry polynomial
\(C(x) = g(x)h(x) - f(x)\) encodes the discrepancy between formal and integer
multiplication. Since \(C(b) = 0\), we factor \(C(x) = (x-b)Q(x)\), and the
companion matrix \(M\) of \(Q\) has spectral properties that approximate Euler
factors of \(\zeta(s)\).

We emphasize: this is an \textbf{approximation} (\(R \neq 1\)), not an exact
identity. The correction is a genuine feature of the carry ensemble, not
removable by a change of averaging.

\subsubsection{Notation}\label{notation}

\(N = pq\) (semiprime), \(g(x), h(x), f(x)\) (digit polynomials),
\(C(x) = g(x)h(x) - f(x)\), \(Q(x) = C(x)/(x-b)\), \(M\) (companion matrix of
\(Q\)), \(D = \deg(Q)\). The symbol \(R(l,s)\) denotes the carry correction
factor defined in §4; this is distinct from the sector ratio \(R(K)\) used in
companion papers {[}E{]}, {[}P1{]}, {[}P2{]}.

\begin{center}\rule{0.5\linewidth}{0.5pt}\end{center}

\newpage

\subsection{2. The Carry Representation
Theorem}\label{the-carry-representation-theorem}

\textbf{Theorem 1} (CRT). \(Q(x) = C(x)/(x-b)\) has integer coefficients
\(q_i = -c_{i+1}\).

\emph{Proof.} Since \(C(b) = g(b)h(b) - f(b) = pq - pq = 0\):
\((x-b) \mid C(x)\). Comparing coefficients of \(C(x) = (x-b)Q(x)\):
\(C_k = q_{k-1} - bq_k\). Induction with \(q_0 = -c_1\) yields
\(q_k = -c_{k+1}\). \(\square\)

Verified with zero exceptions across 185,000 semiprimes (\(\Delta < 10^{-15}\)).

\begin{center}\rule{0.5\linewidth}{0.5pt}\end{center}

\newpage

\subsection{3. The Unit Leading Carry
Theorem}\label{the-unit-leading-carry-theorem}

\textbf{Theorem 2} (ULC). For binary multiplication (\(b = 2\)) of odd primes
with \(C(x) \neq 0\): the leading carry satisfies \(c_{\text{top}} = 1\).

\emph{Proof.} Since \(C \neq 0\) by hypothesis, \(Q = C/(x - b) \neq 0\). Write
\(D = \deg(Q)\), so \(\deg(C) = D + 1\). By the CRT, the leading coefficient of
\(Q\) satisfies \(q_D = -c_{D+1}\); we identify \(c_{\text{top}} = c_{D+1}\) and
establish \(c_{D+1} = 1\) by a two-sided bound.

\textbf{Carries are non-negative integers.} The carry sequence is defined by
\(c_0 = 0\) and

\[c_{k+1} = \left\lfloor \frac{\text{conv}_k + c_k}{2} \right\rfloor, \qquad \text{conv}_k = \sum_{i+j=k} g_i h_j \geq 0.\]

By induction on \(k\): \(c_0 = 0 \in \mathbb{Z}_{\geq 0}\), and if
\(c_k \in \mathbb{Z}_{\geq 0}\) then
\(c_{k+1} = \lfloor(\text{conv}_k + c_k)/2\rfloor \in \mathbb{Z}_{\geq 0}\),
since the floor of a non-negative rational is a non-negative integer. Hence
\(c_k \in \mathbb{Z}_{\geq 0}\) for all \(k\).

Carry termination (\(c_{D+2} = 0\)). Let \(m = \deg(g)\), \(n = \deg(h)\). For
\(k > m + n\): \(\text{conv}_k = 0\) (every term \(g_i h_j\) with \(i + j = k\)
has \(i > m\) or \(j > n\), giving \(g_i h_j = 0\)). Since \(pq < b^{m+n+2}\),
the product has at most \(m + n + 2\) digits, so \(f_k = 0\) for
\(k > m + n + 1\). At any position where \(\text{conv}_k = f_k = 0\), the carry
recurrence reduces to \(c_{k+1} = \lfloor c_k/2 \rfloor\). Since
\(\lfloor a/2 \rfloor < a\) for every positive integer \(a\) (and \(b = 2\)),
the sequence \((c_k)_{k > m+n+1}\) is strictly decreasing while positive and
bounded below by \(0\); it therefore reaches \(0\) in finitely many steps. Fix
any index \(T > m + n + 1\) with \(c_T = 0\).

Now consider positions \(k\) with \(D + 1 < k \leq T\). The polynomial identity
\(C_k = 0\) (valid since \(k > \deg(C) = D + 1\)) gives
\(\text{conv}_k - f_k = 0\), so the carry recurrence
\(\text{conv}_k + c_k = f_k + 2c_{k+1}\) simplifies to \(c_k = 2c_{k+1}\).
Starting from \(c_T = 0\) and applying \(c_k = 2c_{k+1}\) downward:
\(c_{T-1} = 2c_T = 0\), \(c_{T-2} = 2c_{T-1} = 0\), and by induction \(c_k = 0\)
for every \(k \geq D + 2\). In particular, \(c_{D+2} = 0\): no carry propagates
beyond the leading position.

Upper bound (\(c_{D+1} \leq 1\)). At position \(D + 1\):
\(c_{D+2} = \lfloor(\text{conv}_{D+1} + c_{D+1})/2\rfloor = 0\), so
\(\text{conv}_{D+1} + c_{D+1} \leq 1\) (otherwise the floor would be
\(\geq 1\)). Since \(\text{conv}_{D+1} \geq 0\): \(c_{D+1} \leq 1\).

Lower bound (\(c_{D+1} \geq 1\)). The hypothesis \(C \neq 0\) implies
\(Q \neq 0\), so \(D = \deg(Q)\) is well-defined and the leading coefficient
satisfies \(q_D \neq 0\) (by the definition of polynomial degree:
\(Q = q_D x^D + \cdots + q_0\) with \(q_D \neq 0\)). The CRT gives
\(q_D = -c_{D+1}\), hence \(c_{D+1} \neq 0\). Since
\(c_{D+1} \in \mathbb{Z}_{\geq 0}\) (established above), \(c_{D+1} \geq 1\).

Combining: \(c_{D+1} = 1\). \(\square\)

(When \(C = 0\) --- i.e., carry-free multiplication --- \(Q\) is identically
zero, the companion matrix is absent, and the theorem is vacuous.)

Verified 20,000/20,000. For general base \(b\):
\(c_{\text{top}} \in \lbrace{}1, \ldots, b-1\rbrace{}\).

\begin{center}\rule{0.5\linewidth}{0.5pt}\end{center}

\newpage

\subsection{4. The Per-Factor Approximation}\label{the-per-factor-approximation}

\subsubsection{4.1 Statement}\label{statement}

\textbf{Corollary 3} (of CRT). For any semiprime \(N = pq\) in base \(b\) with
\(C(x) \neq 0\), and any \(l\) with \(|l^s| > r_{\max}(Q)\):

\[|\det(I - M/l^s)| = \frac{|Q(l^s)|}{(l^s)^D} = \frac{1}{|1 - b/l^s|} \cdot \frac{|C(l^s)|}{(l^s)^{D+1}}\]

This is an algebraic identity (not statistical), holding for every individual
semiprime. The denominator \(|1 - b/l^s|\) reduces to the Euler factor
\(|1 - l^{-s}|\) when \(b = 1\); for \(b \geq 2\) the ratio
\(|1 - b/l^s|/|1 - l^{-s}| = 1 + O(b/l^{\sigma})\) is absorbed into the
correction factor \(R\).

\emph{Proof.} \(\det(I - M/l^s) = (l^s)^{-D}Q(l^s)\) by the companion matrix
characteristic polynomial. The CRT gives \(C(x) = (x-b)Q(x)\). Since
\(C(b) = 0\), specializing at \(x = l^s\) yields \(Q(l^s) = C(l^s)/(l^s - b)\),
and \(|l^s - b|/(l^s) = |1 - b/l^s|\), giving the stated identity. \(\square\)

The ensemble average defines the correction factor:

\[\langle|\det(I - M_N/l^s)|\rangle_N = \frac{R(l,s)}{|1 - l^{-s}|}\]

\textbf{Observation 4} (empirical). \(R - 1 = O(1/l^{\sigma})\), approximately
universal across primes, with leading term \(c_1/l^{\sigma}\):

\begin{longtable}[]{@{}lll@{}}
\toprule\noalign{}
\(l\) & \(R(l,2) - 1\) & \(l^2(R-1)\) \\
\midrule\noalign{}
\endhead
\bottomrule\noalign{}
\endlastfoot
3 & 0.0186 & 0.167 \\
5 & 0.0068 & 0.170 \\
11 & 0.0014 & 0.167 \\
53 & 0.000060 & 0.168 \\
\end{longtable}

At \(s = 2\): \(l^2(R-1) \approx 0.168\) (cv \textless{} 2\%), approximately
constant across primes. This reflects the finite-\(D\) effective \(c_1\); in the
\(D \to \infty\) limit, \(l^s(R-1)\) is conjectured to converge to
\(c_1 = \pi/18 \approx 0.1745\) (4.3 digits from \(K = 21\) enumeration
{[}E{]}).

\subsubsection{4.2 The Correction Factor}\label{the-correction-factor}

\[\ln R(l,s) = \frac{c_1}{l^s} + \frac{c_2}{l^{2s}} + \frac{c_3}{l^{3s}} + \cdots\]

\begin{itemize}
\tightlist
\item
  \(c_1 = \langle c_{\text{top}-1}\rangle - 1\): the \textbf{trace anomaly},
  universal across \(l\), dependent on \(D\).
\item
  \(c_1 = \pi/18 = 0.17453\ldots\) (conjecture; \({\sim}4.3\) digits from direct
  enumeration at \(K = 21\), consistent with polynomial-geometric convergence
  model \(c_1(K) = \pi/18 + P(K)(1/2)^K\) {[}E{]}). High-statistics Monte Carlo
  rules out the natural candidate \(\ln(2)/4\) at \(> 30\sigma\). The D-even
  component has a closed-form proof (Lemma 1 below); the D-odd component is
  transcendental in \(\lbrace{}\ln 2, \ln 3, 1\rbrace{}\) (PSLQ). The effective
  rate \(\rho(K) = |\Delta(K)/\Delta(K{-}1)|\) approaches \(1/2\) for
  \(K \leq 17\) (\(\rho(17) = 0.505\)), then drops below \(1/2\) at
  \(K = 18\)--\(21\) (\(\rho(21) = 0.333\)), indicating a super-geometric phase
  where \(P(K)\) passes through its maximum magnitude {[}E; experiment B30{]}.
  The base rate \(1/2\) is the Diaconis--Fulman eigenvalue \(\lambda_2 = 1/b\),
  proved for all bases \(b\) {[}experiment B29; G, G04{]}.
\end{itemize}

\textbf{Lemma 1 (D-even trace anomaly component).}

\[P_e \cdot c_1^{(e)} = 1 + 3\ln(3/4),\]

where

\[P_e = P(\text{D-even}) = 2(1-\ln 2), \qquad c_1^{(e)} = E[c_{\mathrm{top}-1} - 1 \mid \text{D-even}].\]

\emph{Proof.} Let \(X, Y \sim U[0,1)\) represent the fractional parts of the two
factors. The product \(W = (1+X)(1+Y)\) satisfies \(W \in [1,4)\); the D-even
condition is \(W \geq 2\). For D-even products the cascade is degenerate {[}E,
§8.1{]}, so

\[c_{\mathrm{top}-1} = \lfloor W \rfloor - 1\]

. Since \(\lfloor W \rfloor \in \lbrace{}2, 3\rbrace{}\) on \([2, 4)\), the only
non-zero contribution to \(c_{\mathrm{top}-1} - 1\) comes from
\(\lfloor W \rfloor = 3\), i.e., \(W \geq 3\). Hence
\(P_e \cdot c_1^{(e)} = P(W \geq 3)\). The integration domain
\(\lbrace{}(x,y) \in [0,1)^2 : (1+x)(1+y) \geq 3\rbrace{}\) has boundary
\(y = (2-x)/(1+x)\), which enters \([0,1)\) at \(x = 1/2\). Therefore

\[P(W \geq 3) = \int_{1/2}^{1}\!\int_{(2-x)/(1+x)}^{1} dy\,dx = \int_{1/2}^{1}\!\left(1 - \frac{2-x}{1+x}\right)dx = \int_{1/2}^{1}\frac{2x-1}{1+x}\,dx\]

\[= \left[2x - 3\ln(1+x)\right]_{1/2}^{1} = (2 - 3\ln 2) - (1 - 3\ln(3/2)) = 1 + 3\ln(3/4). \qquad\square\]

\emph{Remark.} The trace anomaly \(c_1\) of the Euler product correction studied
in this paper is the leading spectral constant of the carry correction factor;
it is related to, but distinct from, the macroscopic cascade trace anomaly
\(\Delta R\) studied in {[}E{]}, which concerns the full sector ratio of the
cascade valuation. - \(c_2 \approx 0.181\), also universal in \(l\). -
Base-specific: \(c_1(b)\) varies with base; \(c_1(3) \approx 0.5985\) (candidate
\(\ln 3 - 1/2\) {[}G{]}), \(c_1(10) \approx 4.20\) {[}E{]}.

\subsubsection{\texorpdfstring{4.3 Structural Properties of
\(R\)}{4.3 Structural Properties of R}}\label{structural-properties-of-r}

\begin{enumerate}
\def\labelenumi{\arabic{enumi}.}
\tightlist
\item
  \textbf{Jensen gap is genuine}: the geometric mean
  \(\exp(\langle\ln|\det|\rangle)\) still yields \(R_{\text{geom}} \neq 1\). The
  Jensen gap accounts for 2--6\% of the total correction, confirming that \(R\)
  is not an averaging artifact.
\item
  \textbf{Character-blind determinant}: the spectral determinant shows no
  sensitivity to Dirichlet character structure --- no direct connection to
  \(L\)-functions emerges at this level.
\item
  \textbf{Base specificity and the Angular Uniqueness Theorem {[}G{]}}:
  \(c_1 = \pi/18\) is specific to base 2, where the D-parity boundary
  \((1+X)(1+Y) = b\) becomes a straight line \(\alpha + \beta = \pi/4\) in
  angular coordinates (\(X = \tan\alpha\), \(Y = \tan\beta\)). This is
  algebraically unique: \(\tan(\alpha+\beta) = (b-1-P)/(1-P)\) is constant iff
  \(b = 2\) {[}G{]}. For \(b > 2\), \(c_1(b)\) does not involve \(\pi\).
\end{enumerate}

\begin{center}\rule{0.5\linewidth}{0.5pt}\end{center}

\newpage

\subsection{5. The Exact Correction Series}\label{the-exact-correction-series}

The CRT yields an exact algebraic identity:

\[h(l) = \langle|\det(I-M/l)|\rangle \cdot (1-1/l) = 1 + \sum_{k=1}^{D-1}\frac{\alpha_k}{l^k}\]

where
\(\alpha_k = \langle c_{\text{top}-k}\rangle - \langle c_{\text{top}-k+1}\rangle\).

Verified to \(\Delta < 10^{-15}\) across 155,000 semiprimes. The coefficients
\(\alpha_k \to 1/4\) (base 2) for \(k \geq 5\) --- consistent with
\(E[c_j] = (j-1)/4\) {[}F{]} --- with boundary deviations at \(k = 1, 2\)
explained by the carry Markov chain's 2-step mixing time {[}A{]}.

\begin{center}\rule{0.5\linewidth}{0.5pt}\end{center}

\newpage

\subsection{\texorpdfstring{6. Convergence Domain at
\(\sigma = 1/2\)}{6. Convergence Domain at \textbackslash sigma = 1/2}}\label{convergence-domain-at-sigma-12}

\subsubsection{6.1 The Spectral Radius
Constraint}\label{the-spectral-radius-constraint}

The bound \(r_{\max} \leq b = 2\) {[}A, Conjecture 10{]} (verified
computationally for 15,000+ semiprimes) ensures convergence of
\(\det(I - M/l^s)\) for \(l^{\sigma} > 2\), i.e., \(l \geq 5\) at
\(\sigma = 1/2\). The weaker Eneström-Kakeya bound \(r_{\max} \leq 3\) (proved)
gives \(l \geq 11\).

\subsubsection{6.2 The Small-Prime Problem}\label{the-small-prime-problem}

For \(l = 2, 3\): the bound \(r_{\max} \leq 2\) does not guarantee convergence
at \(\sigma = 1/2\) (since \(l^{1/2} < 2\)). Additionally, \textasciitilde3\% of
semiprimes have \(r_{\max} > \sqrt{2}\), requiring the renormalization
\(R(l,s)\).

The carry barrier at \(\sigma = 1/2\) is a statement about
\textbf{ensemble-averaged convergence}: the \(\alpha_k\) correction decay rate
\((b-1)/b\) (see {[}A, §5{]}) provides square-root cancellation in carry sums,
but absolute per-factor convergence at small primes requires renormalization.

\begin{center}\rule{0.5\linewidth}{0.5pt}\end{center}

\newpage

\subsection{7. Discussion}\label{discussion}

\subsubsection{7.1 What This Framework
Achieves}\label{what-this-framework-achieves}

\begin{enumerate}
\def\labelenumi{\arabic{enumi}.}
\tightlist
\item
  \textbf{Two proved structural theorems.} The CRT (Theorem 1) and ULC (Theorem
  2) are unconditional, elementary results about base-\(b\) multiplication. They
  hold for every semiprime and require no statistical assumptions.
\item
  \textbf{An exact algebraic decomposition.} Corollary 3 expresses
  \(|\det(I - M/l^s)|\) as a ratio involving the carry-base factor
  \(|1 - b/l^s|^{-1}\), for each individual semiprime. This is a polynomial
  identity, not a statistical observation. The connection to the Euler factor
  \(|1 - l^{-s}|^{-1}\) introduces a correction \(|1 - b/l^s|/|1 - l^{-s}|\)
  that is absorbed into \(R\).
\item
  \textbf{An exact correction series.} The identity
  \(h(l) = 1 + \sum \alpha_k / l^k\) (§5) is proved algebraically from CRT and
  verified to \(\Delta < 10^{-15}\). The coefficients \(\alpha_k\) are
  carry-chain observables computable from first principles.
\item
  \textbf{A falsifiable prediction.} The universal expansion
  \(\ln R = c_1/l^s + \cdots\) (§4.2) predicts that \(c_1\) is base-specific:
  \(c_1(2) \approx 0.1745\), \(c_1(3) \approx 0.5985\),
  \(c_1(10) \approx 4.20\). These are independently verifiable. The Angular
  Uniqueness Theorem {[}G{]} explains why \(\pi\) appears only in base 2.
\item
  \textbf{An honest control experiment.} The isolated correction
  \(\prod R(l,s)\) is smooth with no structure at Riemann zeros {[}H{]},
  demonstrating that the framework does not overclaim.
\end{enumerate}

\subsubsection{7.2 Unconditional vs.~Conjectural
Content}\label{unconditional-vs.-conjectural-content}

Remark (Independence from \(c_1 = \pi/18\)). The structural results of this
paper --- CRT (Theorem 1), ULC (Theorem 2), the algebraic decomposition
(Corollary 3), and the exact correction series (§5) --- are unconditional. They
hold for every base, every semiprime, and any value of \(c_1\). The conjectured
value \(c_1 = \pi/18\) quantifies the leading correction term but is not
required for any theorem. If a future proof established \(c_1 = \kappa\) for
some other constant \(\kappa \in (0,1)\), every structural result would hold
with \(\kappa\) replacing \(\pi/18\). The framework's validity is logically
independent of the specific numerical value of \(c_1\).

\subsubsection{7.3 What It Does Not Achieve}\label{what-it-does-not-achieve}

\begin{itemize}
\tightlist
\item
  \textbf{Exact identity.} \(R \neq 1\) is genuine, not an artefact. The Jensen
  gap test confirms \(R_{\text{geom}} \neq 1\) independently.
\item
  \textbf{Functional equation.} The carry framework provides no analogue of
  \(\xi(s) = \xi(1-s)\). Since the functional equation is what forces Riemann
  zeros onto \(\sigma = 1/2\), no amount of spectral structure in \(R\) can
  substitute for it. The carry-zeta product is not \(\zeta(s)\) and, without a
  functional equation, is not a generalized \(L\)-function.
\item
  \textbf{Analytic continuation.} The carry product exists only for
  \(\sigma > 0\).
\item
  \textbf{Information about Riemann zeros.} Define
  \(Z_{\text{carry}}(s) = \prod_l \langle|\det(I - M_l/l^s)|\rangle\). A control
  experiment {[}H{]} decomposes this as
  \(Z_{\text{carry}} = Z_{\text{Euler}} \cdot \prod R\) and shows that the pure
  partial Euler product \(\prod |1-l^{-s}|^{-2}\) already produces identical
  minima near Riemann zeros (\(\Delta t = 0.02\)--\(0.15\)), while the isolated
  carry correction \(\prod R(l,s)\) is a smooth function with no structure at
  the zeros. The carry framework adds no information about Riemann zeros beyond
  the Euler product itself.
\item
  \textbf{Recovery of prime distribution.} The carry ensemble averages over
  semiprimes \(N = pq\). No mechanism has been identified by which such an
  average could recover information about the distribution of individual primes,
  which is what \(\zeta(s)\) encodes via the von Mangoldt explicit formula.
\end{itemize}

\subsubsection{7.4 Open Questions}\label{open-questions}

\begin{enumerate}
\def\labelenumi{\arabic{enumi}.}
\item
  Is \(c_1 = \pi/18\) exact? The natural candidate \(\ln(2)/4\) is ruled out at
  \(> 30\sigma\) by Monte Carlo {[}E{]}. Direct enumeration at \(K = 21\)
  (\(4^{21} \approx 4.4 \times 10^{12}\) configurations) confirms the
  convergence of \(c_1(K)\) toward \(\pi/18 = 0.17453\ldots\) (matching to 4.3
  significant digits at \(K = 21\)). The residual is consistent with the
  polynomial-geometric model \(c_1(K) = \pi/18 + P(K)(1/2)^K\) {[}E{]};
  model-free evidence is limited to the 4.3-digit agreement. The D-odd
  conditional is transcendental in \(\lbrace{}\ln 2, \ln 3, 1\rbrace{}\)
  (PSLQ-proved). The convergence \(c_1(K) \to \pi/18\) is governed by the
  Diaconis--Fulman eigenvalue \(\lambda_2 = 1/b\), proved for general base \(b\)
  {[}experiment B29; G, G04{]}. At \(K = 18\)--\(21\), the effective rate drops
  below \(1/2\) (super-geometric phase). The appearance of \(\pi\) is explained
  by the Angular Uniqueness Theorem {[}G, Theorem{]}. A formal proof of
  \(c_1 = \pi/18\) remains the central open problem.
\item
  What does \(\prod_l R(l,s)\) converge to? This quantifies the gap between
  carries and \(\zeta(s)\).
\item
  Can a modified ensemble (weighting, different integer class) eliminate \(R\)?
\end{enumerate}

\begin{center}\rule{0.5\linewidth}{0.5pt}\end{center}

\newpage

\subsection{8. Reproducibility}\label{reproducibility}

28 experiment scripts and 1 shared utility module in \texttt{experiments/}. Key
scripts: \texttt{B15\_CRT\_ULC\_proof.py},
\texttt{B14\_c2\_analytical\_decomposition.py},
\texttt{B26\_highprec\_c1\_measurement.py}, \texttt{B27\_jensen\_gap\_test.py},
\texttt{B28\_multibase\_c1.py}.

Requirements: Python 3.8+, NumPy, SciPy.

\begin{center}\rule{0.5\linewidth}{0.5pt}\end{center}

\newpage

\subsection{References}\label{references}

\begin{enumerate}
\def\labelenumi{\arabic{enumi}.}
\tightlist
\item
  P. Diaconis, J. Fulman, ``Carries, Shuffling, and Symmetric Functions,''
  \emph{Adv. Appl. Math.} 43(2), 176--196, 2009.
\item
  J. Holte, ``Carries, Combinatorics, and an Amazing Matrix,'' \emph{Amer. Math.
  Monthly} 104(2), 138--149, 1997.
\item
  D. E. Knuth, \emph{The Art of Computer Programming, Vol. 2}, 3rd ed.,
  Addison-Wesley, 1997.
\item
  E. C. Titchmarsh, \emph{The Theory of the Riemann Zeta-Function}, 2nd ed.,
  Oxford, 1986.
\item
  {[}A{]} Companion paper: ``Spectral Theory of Carries in Positional
  Multiplication,'' this series. doi:10.5281/zenodo.18895593
\item
  {[}E{]} Companion paper: ``The Trace Anomaly of Binary Multiplication,'' this
  series. doi:10.5281/zenodo.18895604
\item
  {[}F{]} Companion paper: ``Exact Covariance Structure of Binary Carry
  Chains,'' this series (\texttt{carry-arithmetic-F-covariance-structure} repo).
  doi:10.5281/zenodo.18895607
\item
  {[}G{]} Companion paper: ``The Angular Uniqueness of Base 2 in Positional
  Multiplication,'' this series. doi:10.5281/zenodo.18895601
\item
  {[}H{]} S. Alimonti, \emph{Carry Polynomials and the Partial Euler Product: A
  Control Experiment}, carry-arithmetic-H-euler-control (2026).
  doi:10.5281/zenodo.18895603
\end{enumerate}

\begin{center}\rule{0.5\linewidth}{0.5pt}\end{center}

\emph{CC BY 4.0. Code: MIT License.}

\end{document}
